\section{真空中的静电场}
\sourceproblem{选择1}{21}{27}
\problemtext 两无限均匀带电平面A、B平行，面密度分别为 $\sigma,2\sigma$，B在右，求两面间及B右侧的有向场强。
\workings 取向右为正，单无限平面两侧场强的法向分量分别为 $\pm\sigma/(2\varepsilon_0)$。叠加得 $E_{\rm between}=-\sigma/(2\varepsilon_0)$，$E_{\rm right}=3\sigma/(2\varepsilon_0)$，选C。
\crosscheck 跨过B时有向场强跳变 $2\sigma/\varepsilon_0$，符合面电荷边界条件。
\sourceproblem{选择2}{21}{27}
\problemtext 边长 $b$ 的正方形中心放置点电荷 $Q$，求顶点场强大小。
\workings 中心至顶点距离 $r=b/\sqrt2$，库仑场强大小为 $|Q|/(4\pi\varepsilon_0r^2)=|Q|/(2\pi\varepsilon_0b^2)$；原选项B采用 $Q>0$ 约定。
\crosscheck 若 $Q<0$，场指向中心，大小仍非负。
\sourceproblem{选择3}{21}{27}
\problemtext 辨别点电荷、无限直线、无限平面及球面外的场强表达式。
\workings 点电荷场 $\bm E=q\bm e_r/(4\pi\varepsilon_0r^2)$，故A正确。无限线电荷场 $\bm E=\lambda\bm r_\perp/(2\pi\varepsilon_0r_\perp^2)$，原B的分母指数不符；单无限带电平面场强大小 $|\sigma|/(2\varepsilon_0)$，原C差一倍；均匀球面外 $\bm E=\sigma R^2\bm e_r/(\varepsilon_0r^2)$，原D对单位矢量的系数误写为 $r^{-3}$。选A。
\crosscheck 区分位矢和径向单位矢量；高斯面分别取柱面、平面薄盒和球面可独立推出各式。
\sourceproblem{选择4}{21}{27}
\problemtext 图2给出内区场强为零、表面外场强随距离平方反比衰减的曲线，辨别电荷分布。
\workings 均匀带电球面内 $Q_{\rm enc}=0$，外部 $E=|Q|/(4\pi\varepsilon_0r^2)$，符合图示，选C。
\crosscheck 无限长圆柱外场衰减为 $1/r$；均匀带电实心球内部场强正比于 $r$。
\sourceproblem{选择5}{22}{28}
\problemtext 图3半径 $R$ 半球曲面处于沿 $x$ 轴的均匀电场，场方向平行于半球底面，求曲面外向通量。
\workings 加上圆形底面构成闭合面，内部无电荷，总通量为零。底面法向垂直于 $\bm E$，其通量为零，故曲面通量亦为零，选A。
\crosscheck 若电场改为垂直底面，曲面通量才为 $\pm E\pi R^2$；不能仅凭曲面积 $2\pi R^2$ 直接乘场强。
\sourceproblem{选择6}{22}{28}
\problemtext 高斯面内净电荷 $8.850\times10^{-12}\,\mathrm C$，判断闭合面电通量。
\workings $\Phi_E=Q_{\rm enc}/\varepsilon_0$，按教材近似 $\varepsilon_0=8.85\times10^{-12}\,\mathrm{F/m}$，得 $\Phi_E\approx1\,\mathrm{N\,m^2/C}$，选C。
\crosscheck 仅总通量确定，不代表各面元通量相同；使用更精确常数时数值约0.9995。
\sourceproblem{选择7}{22}{28}
\problemtext 两同心均匀带电球面半径 $R_a<R_b$，电荷 $Q_a,Q_b$，求 $r>R_b$ 处场强大小。
\workings 由球对称及高斯定理，
\[
\bm E=\frac{Q_a+Q_b}{4\pi\varepsilon_0r^2}\bm e_r,\qquad
E=\frac{|Q_a+Q_b|}{4\pi\varepsilon_0r^2}.
\]
原选项B为净电荷非负时的表达式，或理解为有向径向分量。
\crosscheck 若两球面净电荷抵消，则外场为零；场强大小不可为负。
\sourceproblem{选择8}{22}{28}
\problemtext 根据高斯定理判断闭合面内净电荷与面上各点场强的关系。
\workings $\oint_S\bm E\cdot\mathrm d\bm S=Q_{\rm enc}/\varepsilon_0$。净电荷为零只意味着总通量为零，面上仍可能有非零场强，因此选A。
\crosscheck 外电荷能产生非零场但对闭合面净通量为零；即使面上场处处为零，也只能推知净内电荷为零，不能断言内部处处无电荷。
\sourceproblem{选择9}{22}{28}
\problemtext 根据电势定义辨别电势和试验电荷电势能的关系。
\workings 在选定同一电势零点后，$U=q_0V$，$V=U/q_0$，所以电势等于单位正试验电荷的电势能数值，选D。
\crosscheck 将单位正电荷从该点准静态移至零势点，电场力做功为 $V$ 对应的能量，外力做功为其相反数；区分两类力。
\sourceproblem{选择10}{22}{28}
\problemtext 图4实线为电场线，虚线为等势面，判断A、B、C处场强和电势次序。
\workings 图示电场线箭头指向左侧，沿电场方向电势降低，故 $V_A<V_B<V_C$。按教材电场线密度约定，左侧较密，$E_A>E_B>E_C$，选D。
\crosscheck 场强排序为示意图的定性读图结论；定量值需电势间隔、空间尺度或实际场分布数据。

\sourceproblem{填空1}{23}{29}
\problemtext 根据定义填写静电场强与电场力的关系。
\workings $\bm E=\bm F/q_0$，即单位正试验电荷所受的电场力。
\crosscheck 试验电荷须足够小而不改变原场；负电荷受力方向与场强相反。
\sourceproblem{填空2}{23}{29}
\problemtext 正试验电荷 $4\times10^{-9}\,\mathrm C$ 受到向下 $8\times10^{-9}\,\mathrm N$ 的力，求场强。
\workings $E=F/q_0=2\,\mathrm{N/C}$，方向竖直向下。
\crosscheck 已放大PDF29核实两个指数均为负9；正电荷受力与场同向。
\sourceproblem{填空3}{23}{29}
\problemtext 两平行无限带电面A、B之间场强为 $E_0$ 向左，两外侧场强为 $E_0/3$ 向外，求两面电荷密度。
\workings 取右向为正，三区场强依次为 $-E_0/3,-E_0,+E_0/3$。由法向场跳跃条件，
\[
\sigma_A=-2\varepsilon_0E_0/3,\qquad
\sigma_B=4\varepsilon_0E_0/3.
\]
\crosscheck 两面场叠加：中间为 $(\sigma_A-\sigma_B)/(2\varepsilon_0)=-E_0$，右侧为 $(\sigma_A+\sigma_B)/(2\varepsilon_0)=E_0/3$，与图5一致。
\sourceproblem{填空4}{23}{29}
\problemtext 点电荷 $q$ 的场中，以距电荷 $R$ 的P点为零势点，求距电荷 $r$ 的P'点电势。
\workings $V(r)=q(1/r-1/R)/(4\pi\varepsilon_0)$。
\crosscheck $r=R$ 时为零，$r\to\infty$ 时不一定为零；电势零点由题设指定。
\sourceproblem{填空5}{23}{29}
\problemtext 闭合面内有 $2q$ 与 $-q$，求外向电通量。
\workings 净包围电荷为 $q$，故 $\Phi=q/\varepsilon_0$。
\crosscheck 与两电荷在面内的位置无关；通量可有符号。
\sourceproblem{填空6}{23}{29}
\problemtext 试验电荷 $q$ 在点电荷 $Q$ 的场中沿半径 $R$ 的四分之三圆弧运动，再至无穷远，求电场力功。
\workings 圆弧上电势相同，圆弧段功为零；全程 $W=q[V(R)-V(\infty)]=qQ/(4\pi\varepsilon_0R)$。
\crosscheck 静电场保守，功只依赖两端电势，不依赖圆弧长度。
\sourceproblem{填空7}{23}{29}
\problemtext 正电荷 $1.6\times10^{-19}\,\mathrm C$ 从A至B电场力做功 $3.2\times10^{-15}\,\mathrm J$，求沿其余圆弧返回的功及以B为零势时A的电势。
\workings $W_{BA}=-W_{AB}=-3.2\times10^{-15}\,\mathrm J$；$V_A-V_B=W_{AB}/q=2.0\times10^4\,\mathrm V$。
\crosscheck 闭合路径总功为零，已放大确认两个指数；与往返路径长度无关。
\sourceproblem{填空8}{23}{29}
\problemtext 匀强场 $\bm E=50\bm i+20\bm j$（V/m），A坐标 $(4,2)$、B坐标 $(2,0)$（m），求 $U_{AB}$。
\workings $V_A-V_B=-\bm E\cdot(\bm r_A-\bm r_B)=-50(2)-20(2)=-140\,\mathrm V$。
\crosscheck A在沿电场方向投影较大的一侧，电势更低；反向电势差为正140伏。
\sourceproblem{填空9}{23}{29}
\problemtext 点电荷 $q$ 场中，试验电荷 $q_0$ 从距源 $r_a$ 的a点至距源 $r_b$ 的b点，求外力克服电场力所做功。
\workings 准静态外力功等于势能增加，$W_{\rm ext}=qq_0(1/r_b-1/r_a)/(4\pi\varepsilon_0)$。
\crosscheck 若外力还改变动能，则外力总功须另加 $\Delta K$；这里只求克服电场力的功。
\sourceproblem{计算1}{24}{30}
\problemtext 半径 $R$、总正电荷 $Q$ 的均匀细半圆环位于左半平面，求环心场强。
\workings $\lambda=Q/(\pi R)$。令极角 $\theta\in[\pi/2,3\pi/2]$，$\mathrm d\bm E=-\lambda(\cos\theta\bm i+\sin\theta\bm j)\mathrm d\theta/(4\pi\varepsilon_0R)$。积分得
\[
\bm E=\frac{Q}{2\pi^2\varepsilon_0R^2}\bm i.
\]
\crosscheck 纵向分量对称抵消，水平分量均朝右；量纲为场强。
\sourceproblem{计算2}{24}{30}
\problemtext 半径 $R$、均匀体电荷密度 $\rho$ 的实心球，求内外场强大小。
\workings 由球对称，高斯面包围电荷为 $4\pi\rho\min(r^3,R^3)/3$，故
\[
E(r)=\begin{cases}|\rho|r/(3\varepsilon_0),&0\le r\le R,\\
|\rho|R^3/(3\varepsilon_0r^2),&r>R.\end{cases}
\]
方向由 $\rho$ 符号决定。
\crosscheck 表面连续，中心为零，外场相当于总电荷集中在球心。
\sourceproblem{计算3}{25}{31}
\problemtext 总电荷 $q$ 均匀分布在长 $2l$ 的细杆上，P在延长线上距近端 $a>0$，求电势。
\workings 以杆左端为原点，$x_P=2l+a$，取无穷远零势，
\[
V(P)=\frac{q}{8\pi\varepsilon_0l}\int_0^{2l}\frac{\mathrm dx}{2l+a-x}
=\frac{q}{8\pi\varepsilon_0l}\ln\frac{a+2l}{a}.
\]
\crosscheck 远场 $a\gg l$ 时主项为 $q/(4\pi\varepsilon_0a)$；对数自变量无量纲。
\sourceproblem{计算4}{25}{31}
\problemtext 两同心均匀带电球面半径0.03、0.10米，带等量异号电荷，电势差大小450伏，求内球电荷量。
\workings 内球电荷记 $Q$，外球为 $-Q$，则
\[
V_1-V_2=\frac{Q}{4\pi\varepsilon_0}(1/R_1-1/R_2).
\]
所以 $|Q|=4\pi\varepsilon_0(450)/(1/0.03-1/0.10)\approx2.15\times10^{-9}\,\mathrm C$。
\crosscheck 若明确内球电势更高，则 $Q>0$；只给电势差大小时不能确定符号。外场为零。
